\section{Reviewer 1}

\subsection{Major Comments}

\begin{enumerate}
	\item First, I think the new framing of the paper is very effective. The objective is clearly to make a case for a complex networks approach to understanding civil conflicts, and to provide an application of this approach through the use of AME for the analysis of the Nigerian conflict. I recommend that the authors add a paragraph in the discussion section that addresses any limitations that readers should consider given that the results are from an empirical case study of a single conflict ... Are there limits (e.g., geographic, demographic, temporal) regarding the types of civil conflicts to which the results are generalizable? Should researchers extend the analysis to include more cases in the future?
	\begin{itemize}
		\item \textcolor{blue}{ \emph{
		Insert great response.
		}}
	\end{itemize}	
	\item Second, I recommend mentioning somewhere in the main body of the paper that (1) TERGM has been used to model conflict in the past, and (2) the appendix presents a robustness analysis using TERGM. Also, the authors should note which software was used for the TERGM analysis.
	\begin{itemize}
		\item \textcolor{blue}{ \emph{
		Insert great response.
		}}
	\end{itemize}
	\item Third, it appears that my point on evaluating robustness to `k' was not clearly communicated. I was asking the authors to vary `k' for what it represents in their own notation -- the number of groups into which the data was split for the cross-validation.
	\begin{itemize}
		\item \textcolor{blue}{ \emph{
		Insert great response.
		}}
	\end{itemize}	
\end{enumerate}